\documentclass[11pt]{article}
\usepackage{pwdoc}
\title{Firm-level price wedges for pricewedge.com\\[2pt]\large Three design decisions and the evidence behind them}
\author{Christian Opp}
\date{15 September 2026}
\begin{document}\maketitle

\section*{Summary}

The website reports a price wedge for an individual firm. Rebuilding the
paper's construction from raw CRSP and Compustat data exposed three choices
that the paper does not discuss and that move a firm's number by ten to twenty
percentage points (pp). This note records what was chosen and why.

\begin{enumerate}[nosep,leftmargin=*]
\item \textbf{What the firm-level number means.} It is the price wedge that
firms with the firm's characteristics realized, each firm counting once,
with the value-weighted average of all firm-level wedges held at the
market's wedge of zero. The paper's published firm-level wedges do not have
this property: they run about 18\pp{} too high for the typical firm, because
the mapping describes each portfolio by its average \emph{member} while its
wedge belongs to its \emph{dollars}. Section~2.
\item \textbf{How the characteristics are defined.} As the paper the anomaly
comes from defines them. Six of the 57 characteristics in the paper's data
depart from their sources; five are corrected and one is dropped.
Section~3.
\item \textbf{Which sample.} All available formation dates (July 1964 to
January 2011, cash flows to December 2025), and no alternative samples on
the website. Section~4.
\end{enumerate}

The evidence is in two companion documents: the quality report (the mapping
and the level tests) and the signal log (every characteristic's
construction and its check against the paper's own data). Every number here
is produced by a script in the pipeline and can be regenerated.

\section{The price wedge, and the one calibration that defines it}

The price wedge is the log deviation of a stock's market price from its
fundamental value, the present value of its future dividends under a
benchmark stochastic discount factor (SDF). For a portfolio formed at date
$t$ it is estimated as
\[
\mathrm{PW}_t \;=\; -\log\, \mathbb{E}\!\left[\,\frac{\sum_{s=1}^{180} m_{t,t+s} D_{t+s} + m_{t,t+180} P_{t+180}}{P_t}\,\right],
\]
the expectation being the average over formation dates, one per month. A
negative wedge means the price is below fundamental value. The SDF is the
CAPM-based one of the paper, with one free parameter, the market price of
risk, chosen so that \textbf{the price wedge of the value-weighted market
portfolio is zero}. Every wedge is therefore relative to the average dollar
in the market. This calibration is kept unchanged throughout.

Two facts about the object matter below.

\emph{A portfolio's wedge is a dollar-weighted quantity.} It is the wedge of
one dollar spread across the portfolio's members in proportion to their
market capitalization, so a few large members set it.

\emph{The average firm is far from the average dollar.} With the paper's
price of risk, the value-weighted market has a wedge of zero by
construction; the equal-weighted market, one dollar in every firm, has a
realized wedge of $-24.6$\pp. The average firm is a small firm, and under
this SDF small firms are strongly underpriced. Any firm-level number has to
say which of the two averages it refers to.

\section{Decision 1: what the firm-level price wedge means}

\subsection*{How the paper carries portfolio wedges to firms}

Each of the 570 characteristic-sorted decile portfolios has a price wedge.
Each also has a set of \emph{portfolio-level characteristics}: the
rank-normalized characteristics of its members, averaged over the members
and over time. A regression of the wedges on the portfolio-level
characteristics gives a rule from characteristics to wedge, and evaluating
that rule at an individual firm's own characteristics gives the firm-level
price wedge. The paper uses three principal components of the
characteristics of the 114 extreme deciles; the pipeline uses a penalized
regression on all 570, which fits better on every portfolio-level test
(quality report, Section~1). That difference is not the decision here.

The decision is hidden in the word \emph{average}. A portfolio's average
characteristics can be taken over its members with equal weights, which
describes its typical member, or with market-capitalization weights, which
describes its typical dollar. For the value-weighted book-to-market decile
10, the typical member has a size rank near 0.4 and the typical dollar a
size rank near 0.98. The paper's code uses equal weights (with the comment
``SHOULD THIS BE VALUE WEIGHTED AVERAGE?????'' beside it); the pipeline had
used capitalization weights. The choice moves every firm-level wedge by ten
to twenty-seven points, and it is the whole reason the pipeline's firm-level
wedges first sat eleven points below the published ones: with equal weights
the published firm-level wedges are reproduced to 1.2\pp{} (slope 1.04,
correlation 0.99).

\subsection*{The test}

Two firm-level series can correlate at 0.95 and differ by 20\pp{} in level,
and the level is what the website prints. The test used here therefore
compares assigned wedges with realized ones. Take a candidate firm-level
wedge; sort firms into deciles on it every month, as one would on a
characteristic; hold each decile for fifteen years, equal-weighted and
value-weighted; compute the decile's realized price wedge with the same SDF;
and regress the realized wedge on the average wedge the candidate had
assigned to the decile's members. A candidate whose levels can be trusted
gives a slope of one and an intercept of zero. The equal-weighted deciles
answer the website's question, whether firms realize what they are told; the
value-weighted deciles check that the candidate respects the calibration.
The calibration gives one more check: the value-weighted average of all
firm-level wedges must equal the market's wedge, zero.

\subsection*{Three candidates}

\begin{tbl}{Three candidates for the firm-level price wedge, and what the level test says about each.}
\begin{tabular}{>{\raggedright\arraybackslash}p{3.1cm}>{\raggedright\arraybackslash}p{5.6cm}rrrr}\toprule
 & & \multicolumn{2}{c}{equal-weighted deciles} & \multicolumn{2}{c}{value-weighted}\\
\cmidrule(lr){3-4}\cmidrule(lr){5-6}
Candidate & What the number claims & slope & intercept & slope & VW mean\\\midrule
1. Published construction & ``the number the paper's method assigns'' & 0.64 & $-18$ & 0.95 & $+9$\\[2pt]
2. Per dollar & ``a dollar in firms like this one is mispriced by $X$'' & 0.67 & $-3$ & 0.79 & $+1$\\[2pt]
3. Per firm, aggregate pinned & ``firms like this one are mispriced by $X$ on average; dollar-weighted, mispricing is zero'' & 0.86 & $-1$ & 0.76 & $0$\\
\bottomrule
\end{tabular}
\tablenote{Candidate 1: value-weighted portfolio wedges regressed on equal-weighted portfolio characteristics. Candidate 2: value-weighted wedges and characteristics. Candidate 3: value- and equal-weighted portfolio sets together, each with characteristics of its own weighting. Intercepts in pp. ``VW mean'' is the value-weighted average of the candidate's firm-level wedges, which should be near zero.}
\end{tbl}

\textbf{Candidate 1} learns ``a portfolio whose typical member is a small
value firm has wedge $W$'' when $W$ was earned by the portfolio's large
members, and then hands the small firm the large firm's wedge. Firms it
calls fairly priced realize about $-25$\pp; only the most underpriced decile
realizes what it is assigned.

\textbf{Candidate 2} is internally consistent but is evaluated far outside
the points it was fitted on: almost every portfolio's dollar-average size
rank is above 0.85, and the website asks about firms at 0.4. The level is
right on average and the extremes are exaggerated (the most underpriced
decile is assigned $-70$\pp{} and realizes $-42$).

\textbf{Candidate 3} adds the same 570 deciles held with equal weights, whose
wedge is the wedge of the average member, each described by the
equal-weighted average of its members' characteristics. A small firm's
number is then interpolated from portfolios whose members resemble it, and
the value-weighted portfolios keep the dollar-weighted total anchored. Out
of sample (mapping fitted on formation dates before 1988, evaluated after)
the slopes hold at 0.87 to 0.99.

\subsection*{Decision}

Candidate 3 is adopted. It is the only candidate whose levels were checked
against realized outcomes at the firm level and found right, equal- and
value-weighted at once; its claim is the one a reader looking up a single
company expects; it keeps the paper's calibration intact; and it changes
rankings very little (all candidates rank firms alike, correlations 0.9 and
above). The published construction is kept as a labelled reference series.

\emph{A consequence for the paper.} The published firm-level wedges are
candidate~1, so any use of them in levels inherits the $-18$\pp{} intercept;
the investment-$q$ regressions use them as a regressor in levels. Rankings
are unaffected.

\section{Decision 2: how the characteristics are defined}

The mapping learns from the wedges of the 570 portfolios, and a portfolio is
defined by how its sorting characteristic was built. Rebuilding the 57
characteristics and checking them against the paper's own firm-month panels
verified 56 of them and found six places where what the paper's data does
differs from the source the paper cites. The rule adopted: \textbf{each
characteristic is defined as the paper the anomaly comes from defines it},
read through Freyberger, Neuhierl and Weber (2020), whose list the paper
takes its characteristics from. The point of the anomalies is to use the
existing specifications, not new ones that arose from an implementation
slip.

\begin{tbl}{The six characteristics whose construction in the paper's data departs from its source.}
\begin{tabular}{>{\raggedright\arraybackslash}p{1.7cm}>{\raggedright\arraybackslash}p{4.4cm}>{\raggedright\arraybackslash}p{4.2cm}>{\raggedright\arraybackslash}p{3.9cm}}\toprule
 & Source & Paper's data & Website\\\midrule
OA, AOA & Sloan (1996): accruals \emph{net of} depreciation & depreciation added & Sloan's sign, lagged assets as in FNW\\[2pt]
SUV & Garfinkel (2009): a daily quantity, the day's volume minus a regression fit from a prior window, over the fit's residual s.d.; FNW: fit on the previous month & the last trading day of a regression run within the same month & daily values from the previous month's fit, averaged over the month\\[2pt]
DTO & Garfinkel (2009): daily turnover minus market turnover, minus a control-period mean; FNW: minus the firm's 180-day median, Nasdaq volume scaled (Anderson and Dyl 2005) & last day's raw turnover minus its 180-day mean; no market adjustment, no scaling & FNW's definition, averaged over the month\\[2pt]
TNOVR & Datar, Naik, Radcliffe (1998); FNW: last month's volume over shares & last month's volume over shares & unchanged (the data follows its source)\\[2pt]
aBEME, aSIZE, aSAT & Asness, Porter, Stevens (2000): 48 Fama-French industries by SIC, equal-weighted industry mean, three firms or more; FNW apply the same to PM and SAT & 27 industry groups of unknown origin & Fama-French 48, equal-weighted mean\\[2pt]
aPM & Soliman (2008) defines PM; the adjustment is FNW's & industry means not reconstructible; long-short alpha zero & dropped\\
\bottomrule
\end{tabular}
\end{tbl}

Two remarks. For SUV and DTO neither source singles out one trading day, so
the paper's last-day values are a construction artifact (the signature of a
daily series collapsed to monthly by keeping the last observation); the
monthly average is Garfinkel's own aggregation over his event window.
Soliman's published paper does not industry-adjust; the adjustment of PM and
SAT is FNW's, built like Asness, Porter and Stevens's, and that is what is
used. Implementation conventions that the sources leave open and that do
not change what a characteristic measures (the 249-day window of the daily
beta, zero-volume days excluded from the volume dispersion, the day-level
bid-ask spread, dividends in dollars over market capitalization) are kept
and documented in \texttt{signals.pdf}.

\section{Decision 3: which formation dates}

The wedges, and the price of risk with them, are averages over formation
dates. Three candidates: the paper's window (July 1964 to December 2002,
cash flows to 2017), all available dates (to January 2011, cash flows to
2025), or a later window.

Extending the paper's window to all available dates changes almost nothing:
the two cross-sections of 560 portfolio wedges correlate at 0.99, one
long-short spread changes sign, firm-level wedges move by 2\pp. Dropping the
early dates is different. The first and second halves of the sample
(1964--87 and 1988--2011) correlate at 0.45, where two halves of a constant
cross-section would correlate at about 0.68 given their estimation noise;
the value ratios reverse sign between halves, profitability spreads shrink,
momentum grows. Any window starting in 1980 or later rests on one or two
fifteen-year cycles (the price of risk swings from 5.6 to 2.5 to 2.8 across
thirds of the sample) and moves the typical firm's wedge by 10 to 20\pp.
Firm rankings are robust to the choice; levels are not.

\textbf{Decision.} All available formation dates. A later window is not a
better estimate of the same quantity but an estimate of a different regime.
The half-sample estimates stay as an internal check on every build; the
website offers no sample choices, because that would present regime
dependence as if it were an alternative estimate.

\section{What the website reports, and what every build must pass}

\textbf{Reported per firm:} the candidate-3 price wedge on all available
dates with the characteristics of Section~3, the firm's percentile among all
firms that month, and a flag when the firm's characteristics lie outside the
range spanned by the portfolios (where the mapping extrapolates). The
published construction is available as a labelled reference. Uncertainty is
stated once, on the method page: about 30\pp{} common to all firms (the
market's realized path), about 3 to 4\pp{} specific to a firm, and about
$\pm12$\pp{} of sensitivity to the half of the sample used.

\textbf{Acceptance tests, all on the latest build:}
\begin{enumerate}[nosep,leftmargin=*]
\item Replication: with the paper's own definitions, at least 56 of 57
characteristics agree with the paper's panels and Table~1 within a point on
both legs.
\item Levels: on equal-weighted deciles of the assigned wedge, slope at
least 0.85 and intercept within 3\pp; on value-weighted deciles, slope at
least 0.75; the decile table reported.
\item Aggregate: the value-weighted average of firm-level wedges within a
point or two of zero.
\item Out-of-sample fit of the mapping (leave one family of characteristics
out) reported and not lower than the previous build's by more than noise.
\item Subsequent returns of the assigned-wedge deciles monotone at five
years.
\item Half-sample check: the mappings fitted on 1964--87 and on 1988--2011
and their firm-level difference from the headline reported.
\end{enumerate}

\section*{Appendix: scripts and sources}

\begin{tbl}{Where each piece lives.}
\begin{tabular}{p{5.6cm}p{8.6cm}}\toprule
\texttt{build\_panel.py -{}-spec site} & the characteristics as their sources define them, to 2025\\
\texttt{portfolio\_wedges.py}, \texttt{portfolio\_profiles.py} & sorts, calibration, wedges and portfolio-level characteristics, value- and equal-weighted (\texttt{-{}-weighting equal})\\
\texttt{firm\_wedges.py -{}-mapping stacked} & the candidate-3 mapping\\
\texttt{realised\_levels.py} & the level test and the aggregate check\\
\texttt{stability.py} & wedges on subsets of formation dates\\
\texttt{validate.py}, \texttt{compare\_panels.py} & replication against Table~1 and the paper's panels\\
quality report, signal log & the evidence (\texttt{doc/quality.pdf}, \texttt{doc/signals.pdf})\\
Literature folder & the source papers (\texttt{Website/Literature/}, indexed in \texttt{\_index.md})\\
\bottomrule
\end{tabular}
\end{tbl}

\end{document}
